% META: {"id": "TEXP-MAT-EX-028", "chapitre": "TEXP-MATRICES-MARKOV", "type_objet": "exercice", "capacites": ["TEXP-MAT-C7"], "capacites_codes": ["C7"], "methodes": ["M4"], "parcours": 2, "competences": ["communiquer","calculer"], "duree_min": 12, "mode_creation": "ex_nihilo", "sources_inspiration": [], "fichier_tex": "chapitres/TEXP-MATRICES-MARKOV/exercices/TEXP-MAT-EX-028.tex", "corrige_tex": "chapitres/TEXP-MATRICES-MARKOV/corriges/TEXP-MAT-CO-028.tex", "status": "approved"}
% BEGIN-VERIFY
% from sympy import *
% x = symbols('x')
% f = x**3 - 3*x
% fp = diff(f, x)
% fpp = diff(f, x, 2)
% assert expand(fp) == 3*x**2 - 3
% assert expand(fpp) == 6*x
% assert solve(fp, x) == [-1, 1]
% assert solve(fpp, x) == [0]
% assert f.subs(x, 0) == 0
% assert f.subs(x, -1) == 2
% assert f.subs(x, 1) == -2
% END-VERIFY
\begin{exercice}{TEXP-MAT-EX-028}{1}{12}
Soit $f(x) = x^3 - 3x$.
\begin{enumerate}
  \item Dresser les tableaux de signes de $f'$ et $f''$.
  \item En deduire les intervalles de convexite et les points d'inflexion.
  \item Esquisser la courbe en placant les points remarquables.
\end{enumerate}

\end{exercice}
